\section{Introduction}
An upstream renewal desk can know an accepted obligation that a downstream execution gateway does not retain. When only a small service symbol crosses that boundary, the provider must choose both its terminal instruction alphabet and the branch-dependent encoder. This is a quantizer-design problem, but its objective and feasible randomization are generated by a contract: an aggregate continuation promise, local participation limits, and the cost of intermediate rather than terminal service. We ask how to choose the alphabet and allocate the promise jointly when individual lottery realizations have an overrun limit and enabling an additional execution level has a cost.

The main result gives a single design formulation for every feasible aggregate promise, from zero to saturation, with continuously chosen execution levels, bounded realization exposure, and charges on the actual selected levels. A branch-specific eligibility prefix and a canonical adjacent lottery reduce each fixed-book problem to concave resource allocation. In the quadratic model, eliminating lottery probabilities then produces finitely many closed quadratic optimization cells. Exact stationary-face enumeration solves the continuous outer problem, including collisions and nonsaturated pathwise participation. This global reference algorithm is exponential, and we state that cost explicitly.

A second result supplies a computational alternative with a uniform guarantee. We round codewords downward, retain every branch promise, and reconstruct the lottery rather than retaining its probabilities. The resulting catalog policy preserves the promise and the risk bound exactly, even at a binding pathwise constraint. Its net-payoff loss has an explicit bound proportional to the catalog spacing, uniformly over promises and risk tolerances. Exact catalog allocation therefore gives certified lower and upper values for the continuous problem, and an additive approximation scheme when the symbol budget is fixed. This connection is what makes a numerical catalog a controlled approximation rather than a change of model.

The saturated contract remains a particularly tractable and informative member of this theory. Deterministic execution yields an ordered one-sided quantizer. Under expected participation, a quadratic optimum has cap-anchored nonhighest levels and one continuously optimized highest level, which can lie strictly between caps because intermediate service meets the shortfall. A crossing-difference identity preserves Monge order after continuous tail minimization. Classical matrix search then gives linear work per dynamic-programming layer. We retain these results, their complete proofs, the all-promise exact-alphabet phase diagram and stability bounds, and the faster saturated charged-catalog recurrence. The general algorithm complements rather than replaces these sharper special-case results.

Ordered quantizer dynamic programming, optimized reproduction points, unbiased random splitting, and matrix search are established ideas \citep{Wu1991,PagesWilbertz2012,CrociEtAl2022,JourdainPages2021}. Likewise, a cost per enabled representative and quadratic active-set analysis are not generic methodological inventions of this paper. The contractual contribution is the interaction of endogenous promises, heterogeneous shortfall service, branch-specific realization ceilings, and a shared selected alphabet. That interaction yields the eligibility-prefix reduction, the exact closed-cell formulation, and a one-sided repair that needs no promise or risk slack. The heterogeneous tail and its optimized Monge cancellation remain a separate structural advantage at saturation.

The operational interpretation is a risk-neutral or actuarial acceptance rule followed by a restricted execution interface; hard realization limits represent a different acceptance institution. We spell out those assumptions and distinguish documented interface-security analogies from evidence of an actual deployed contract. Exact computations map institutional gaps and joint promise--risk--price choices, compare continuous and catalog optima, and time the new algorithms separately from inherited benchmarks. These are synthetic design studies and implementation checks, not customer data or substitutes for mathematical proofs.
